\subsection{The exact reduction for a contracted horizontal section}

The statements C.2.CR.1--C.2.CR.6 below retain the
proved local and global reductions. The later
C.2.UN.7 proves that the contracted irrational
configuration has $e(S)\ge2$ and $\pi_1(S)=1$.
Together with the CDP transgression, or the zero
Euler conclusion C.2.MH.3, this excludes the
configuration entirely under the CDP hypotheses.
The unrestricted genus assertion is separately
proved in C.2.UN.10.

\begin{statement}{C.2.CR.1}{Proved}{Topology and coefficients in the contracted-section case}
Let $Y$ be a smooth compact complex threefold with
$b_1(Y)=b_2(Y)=0$, and let $f:Y\to\PP^1$ be a proper
holomorphic surjection with connected fibres and smooth
general fibre a complex two-torus. Let $mS$
be a full fibre with irreducible nonnormal reduced support, and
write $h=\eta\pi:Z\to T\to S$ for its normalization and minimal
normalization resolution. Suppose that $p:Z\to B$ is an
irrational ruling, that the effective divisor in A.5.3 is
$D=2C+V$, where $C$ is a section and $V$ is vertical, and
that $\pi$ contracts $C$. Put $g=g(B)$, let $\Gamma$ be the
conductor-image scheme, and let $n$ be the number of irreducible
components of $\Gamma_{\rm red}$.

Then the following are exact consequences of these hypotheses:
\[
\begin{gathered}
 H^1(T;\mathbb R)=H^1(T,\OO_T)=H^2(T,\OO_T)=0,
 \qquad \operatorname{length}R^1\pi_*\OO_Z=g,\\
 \pi_1(S)\simeq\pi_1(\Gamma_{\rm red})\simeq\mathbb Z,
 \qquad e(\Gamma_{\rm red})=n,\\
 \chi(\OO_\Gamma)=-1,\qquad
 \chi(\OO_{\Gamma_{\rm red}})\in\{-1,0\},\qquad
 \sum_{x\in\Gamma_{\rm red}}\bigl(\delta_x-r_x+1\bigr)\le1.
\end{gathered}
\tag{C.CR.1}
\]
Here $r_x$ is the number of normalization branches at $x$.
The full inverse image $D_{\rm red}$ is a smooth genus-$g$
section with trees of smooth rational curves attached; all
its intersections are ordinary double points. For adjacent
components of coefficients $a,b$ one has
\[
                         \gcd(a-1,b-1)=1.          \tag{C.CR.2}
\]
Choose a specified sequence of ruling blowdowns to a bundle
$Z_0\to B$. Write its pushed-forward divisor as
$2C_0+\sum_{j=1}^{v}a_jF_j$, with distinct fibres $F_j$
and $a_j>0$, and let $N_0$ count the reverse blowups whose
centre has multiplicity one in the current divisor. Then
\[
                     e(S)=2-2g-v+N_0+n.           \tag{C.CR.3}
\]
These necessary equations are retained in the Euler
calculation. C.2.UN.7 later excludes their simultaneous
geometric occurrence under the stated CDP hypotheses.
\end{statement}
\begin{proof}
The finite-duality calculation (A.RT.5), whose proof does
not assume rationality, gives $H^2(T,\OO_T)=0$. The
degree-one Leray sequence is consequently
\[
 0\longrightarrow H^1(T,\OO_T)\longrightarrow H^1(Z,\OO_Z)
 \longrightarrow H^0(T,R^1\pi_*\OO_Z)\longrightarrow0.
\tag{C.CR.4}
\]
The middle vector space has dimension $g$. At the point
to which $C$ is contracted, formal functions and the
successive infinitesimal exceptional fibres give a
surjection from the completed stalk of $R^1\pi_*\OO_Z$
onto the $H^1$ of the scheme-theoretic exceptional fibre.
Indeed the restriction maps on $H^1$ of these thickenings
are surjective, since their kernels are supported in
dimension at most one. The completed stalk equals the
stalk, because this higher direct image has finite support.
Restriction from the exceptional fibre to $C$ is likewise
surjective on $H^1$. Thus its length is at least
$h^1(C,\OO_C)=g$, whereas (C.CR.4) bounds the total length
by $g$. Equality proves the coherent assertions in the
first line of (C.CR.1).

Connectedness of resolution fibres makes the constant-sheaf
Leray map $H^1(T;\mathbb R)\to H^1(Z;\mathbb R)$ injective.
Restriction from $Z$ to its section $C$ is an isomorphism
on $H^1$, because a ruled bundle has simply connected
fibres and point blowups do not change $H^1$. The image of
$H^1(T;\mathbb R)$ restricts to zero on the contracted
curve $C$. This proves its vanishing.

We describe the actual divisor comparisons. At every ruling
blowdown $\sigma:W\to W'$ with exceptional curve $E$, let
$a$ be the coefficient of $E$ and let $d$ be the multiplicity
of the pushed-forward boundary at the centre. The identity
\[
 K_W+D_W=\sigma^*(K_{W'}+D_{W'})+(1+a-d)E
\tag{C.CR.5}
\]
and numerical triviality give $a=d-1$. The minimal model
boundary is therefore $2C_0+\sum a_jF_j$. Its reduced
support is the section with disjoint rational fibre leaves.
Every reverse blowup preserves the asserted simple-normal-
crossings tree structure: a positive exceptional coefficient
adds a leaf at a smooth boundary point or subdivides an edge
at a boundary node; a zero exceptional coefficient leaves
the reduced boundary isomorphic to its predecessor. A centre
off the current boundary would create coefficient $-1$,
which persists on its strict transform and contradicts the
effectiveness of the original divisor. In particular
$D_{\rm red}$ is connected, and its fundamental group maps
isomorphically to that of $Z$ by inclusion, with both
identified with $\pi_1(B)$ through the section.

This also proves (C.CR.2) without discarding coefficients.
Initially every edge has coefficients $(2,a_j)$ and hence
coprime shifted coefficients. A node blowup inserts the
coefficient $a+b-1$, so its shifted coefficient is
$(a-1)+(b-1)$; the two new edge gcds equal the old gcd.
A positive-coefficient blowup at a smooth point of a
coefficient-$a$ boundary creates coefficient $a-1$ and
an edge with shifted coefficients $a-1,a-2$, whose gcd
is one. These are all the possible operations.

On a minimal normalization resolution, a coefficient-zero
exceptional curve $E$ has $K_Z\cdot E\ge0$. Since
$(K_Z+D)\cdot E=0$, it follows that $D\cdot E=0$.
Such a curve cannot meet $D_{\rm red}$. Connectedness of
exceptional fibres and of $D_{\rm red}$ now gives the same
support equality as in A.5.16:
\[
 D_{\rm red}=h^{-1}(\Gamma_{\rm red}).
\tag{C.CR.6}
\]
Any exceptional fibre outside this set would correspond to
an isolated normal hypersurface singularity of $S$, because
the normal conductor ideal is reflexive and has no isolated
support. A.5.17 excludes every such point in the actual
torus smoothing. Thus $h$ is an isomorphism off the two
closed sets in (C.CR.6).

It follows that $S$ is the topological pushout of
$D_{\rm red}\hookrightarrow Z$ and
$D_{\rm red}\to\Gamma_{\rm red}$. This is the actual
pushout: its map to $S$ is a continuous bijection from a
compact space to a Hausdorff space. Triangulation of the
pair makes the first map a cofibration. Van Kampen applied
to its double mapping cylinder therefore gives
\[
 \pi_1(S)=\pi_1(Z)*_{\pi_1(D_{\rm red})}
                         \pi_1(\Gamma_{\rm red})
          \simeq\pi_1(\Gamma_{\rm red}).
\tag{C.CR.7}
\]
Indeed the first amalgamating map is an isomorphism, and
the second map is trivial, since the section $C$ generates
$\pi_1(D_{\rm red})$ and $h(C)$ is a point. Every component
of $\Gamma_{\rm red}$ is dominated by a vertical rational
curve, so its normalization is $\PP^1$. Its fundamental
group is the free group of rank $b_1(\Gamma_{\rm red})$.

For the exact local comparison with the nearby torus, take
a triangulated regular neighbourhood retracting onto $S$
and put a sufficiently small full fibre tube inside it.
Perturbing loops off the divisor gives a surjection from
the fundamental group of its punctured torus bundle onto
that of $S$. The former is $\mathbb Z^4\rtimes\mathbb Z$,
which is metabelian. A free group of rank at least two
is not metabelian, since it surjects onto $S_4$.
Consequently the rank in (C.CR.7) is zero or one. The
nonzero integral class on every support supplied by D.1.9
excludes zero. This proves the fundamental-group assertion.

Normalization of the rational image components gives the
exact identities
\[
 b_1(\Gamma_{\rm red})
   =\sum_x(r_x-1)-n+1=1,\qquad
 e(\Gamma_{\rm red})=2n-\sum_x(r_x-1)=n.
\tag{C.CR.8}
\]
Put $\epsilon=h^0(S,\omega_S)\in\{0,1\}$.
The vanishings already proved, the normalization exact
sequence and conductor duality give
\[
 h^0(\eta_*\OO_T/\OO_S)=1+\epsilon,\qquad
 h^1(\eta_*\OO_T/\OO_S)=\epsilon,
 \qquad\chi(\OO_\Gamma)=-1.
\tag{C.CR.9}
\]
The nilradical sequence and curve duality give
$\chi(\OO_{\Gamma_{\rm red}})\ge-1$ exactly as in
C.2.VII.2: for $\epsilon=0$ the relevant twisted $H^1$
has dimension at most one, and for $\epsilon=1$ the twist
is trivial and connectedness supplies $h^0=1$.
On a connected curve with rational normalizations,
$\chi(\OO_{\Gamma_{\rm red}})=n-\sum_x\delta_x$.
Together with (C.CR.8) and $\delta_x\ge r_x-1$, this
proves the remaining assertions of (C.CR.1).

Finally $e(Z_0)=4-4g$ and
$e((D_0)_{\rm red})=2-2g+v$. A reverse blowup of
multiplicity at least two increases the Euler numbers
of both the surface and the reduced boundary by one,
leaving their difference unchanged. A multiplicity-one
centre leaves the boundary isomorphic to its predecessor
and increases the complement Euler number by one.
Therefore
\[
 e(Z\setminus D_{\rm red})=2-2g-v+N_0.
\]
Equation (C.CR.6), the excluded off-conductor exceptionals,
and additivity of Euler number now give (C.CR.3).
\end{proof}


\begin{statement}{C.2.CR.2}{Proved}{The ruling genus is at most two}
In the situation of C.2.CR.1, the ruling genus is $g=1$ or
$g=2$. Let $p=\pi(C)$ and $q=\eta(p)$. When $g=2$, the
normalization has just the one point $p$ over $q$, and the
full conductor-image scheme $\Gamma$ is smooth and reduced
at $q$. Its component through $q$ has generic transversal
either an ordinary node or an ordinary cusp.
\end{statement}
\begin{proof}
Work over a small neighbourhood of $q$ and fix a generator
of the invertible dualizing sheaf $\omega_S$. Finite
duality then identifies the conductor on the normal branch
$(T,p)$ with its canonical module:
\[
 I_p=\omega_{T,p}\subset\OO_{T,p}.
\]
The conductor ideal on $S$ is the direct sum of these
normalization-branch ideals; thus $I_p$ itself embeds
as an ideal of $R=\OO_{S,q}$ by putting zero on the other
normalization branches.

Put $J=(\pi_*\omega_Z)_p$. Proper duality and
Grauert--Riemenschneider vanishing give
\[
 \operatorname{length}(I_p/J)
   =\operatorname{length}(R^1\pi_*\OO_Z)_p=g.
 \tag{C.CR.10}
\]
These are the analytic duality and vanishing maps described
in \cite[Sections~2.2.2--2.2.5]{vanStraten1987}.
For precision, the defect of $\OO_T\to R\pi_*\OO_Z$ is
the finite-length module $R^1\pi_*\OO_Z$ in degree one;
duality turns it into the quotient
$\omega_T/\pi_*\omega_Z$, since the positive direct
images of $\omega_Z$ vanish. Finite-length duality
preserves length.

The residue sequence for the smooth exceptional section is
\[
 0\longrightarrow\omega_Z\longrightarrow\omega_Z(C)
   \longrightarrow\omega_C\longrightarrow0 .
 \tag{C.CR.11}
\]
The same vanishing theorem makes its pushforward quotient
at $p$ equal to $H^0(C,\omega_C)$, of dimension $g$.
The sheaf $\pi_*\omega_Z(C)$ embeds in $\omega_T$:
away from the exceptional set its sections are holomorphic
two-forms, and the canonical module of a normal surface
extends these across its isolated singular points.
Comparing lengths with (C.CR.10) proves
\[
 I_p=(\pi_*\omega_Z(C))_p,\qquad
 I_p/J\simeq H^0(C,\omega_C).
 \tag{C.CR.12}
\]
In the fixed dualizing trivialization the exact comparison
of A.5.3 reads
$\omega_Z=\OO_Z(-2C-V)$, so that
$\omega_Z(C)=\OO_Z(-C-V)$. The quotient in (C.CR.12)
is consequently the leading coefficient of order one
along $C$, with the vertical vanishing conditions retained.

Let $\mathfrak m$ be the maximal ideal of $R$. Every
element of $\mathfrak m$ pulls back to a function vanishing
on $C$, because $h(C)=q$. Its square therefore vanishes
to order at least two on $C$. Formula (C.CR.12), including
all vertical conditions in $V$, implies
\[
                 I_p\cap\mathfrak m^2\subset J.
 \tag{C.CR.13}
\]
Indeed an element of $I_p$ is already a section of
$\OO_Z(-C-V)$; an additional zero along $C$ puts it
in $\OO_Z(-2C-V)$.

It follows that
\[
 g=\dim(I_p/J)
 \le \dim\bigl(I_p/(I_p\cap\mathfrak m^2)\bigr)
 \le \dim(\mathfrak m/\mathfrak m^2)\le3.
 \tag{C.CR.14}
\]
If the middle dimension were three, the ideal $I_p$
would generate the whole cotangent space. Nakayama's
lemma applied to $\mathfrak m/I_p$ would give
$I_p=\mathfrak m$. The full conductor ideal, which
contains $I_p$, would then be $\mathfrak m$, making its
image zero-dimensional. This contradicts the pure
one-dimensional Cohen--Macaulay conductor image.
Thus the middle dimension is at most two, proving
$g\le2$.

When $g=2$, choose $x,y\in I_p$ whose linear parts
are independent, and complete them to coordinates
$(x,y,z)$ on the smooth ambient threefold. The full
conductor ideal contains $(x,y)$ in $R$. The defining
equation of $S$ belongs to $(x,y)$, since the conductor
has dimension one. The quotient by $(x,y)$ is therefore
$\mathbb C\{z\}$. A strictly larger proper conductor
ideal would give a zero-dimensional quotient of this
ring. Purity hence forces the full conductor ideal
to be exactly $(x,y)$, proving smoothness and
reducedness of $\Gamma$ at $q$.

There is no other normalization point over $q$.
The two elements $x,y$ were chosen in the single
summand $I_p$, so both vanish identically on any other
local normalization branch. Such a two-dimensional
hypersurface branch cannot be contained in the
one-dimensional smooth germ $\{x=y=0\}$.

Finally conductor duality gives
$Q\simeq\omega_\Gamma\otimes\omega_S^{-1}|_\Gamma$.
Near the smooth reduced curve $\Gamma$ this is a
line bundle, so the transverse normalization quotient
has length one. A reduced plane curve of delta
invariant one is a node or a cusp. One can see this
directly: multiple branches require exactly two
smooth branches with intersection multiplicity one;
a single branch has normalization
$\mathbb C\{u\}$ and a colength-one subring, necessarily
$\mathbb C\{u^2,u^3\}$.
\end{proof}

\begin{statement}{C.2.CR.3}{Proved}{The genus-two node alternative is excluded}
In the genus-two situation of C.2.CR.2, the component
of $\Gamma$ through $q=h(C)$ cannot have generic
transversal an ordinary node. Consequently the
remaining genus-two alternative has an ordinary cusp
at its generic point, and exactly one strict conductor
curve of coefficient two and degree one lies above
that component.
\end{statement}
\begin{proof}
We first record an additional consequence of the
specified boundary blowups in C.2.CR.1: every
coefficient-one component of $D_{\rm red}$ is a leaf.
This holds on the initial section-and-fibres star.
A smooth boundary blowup creates coefficient one
only when the old coefficient is two, and creates
it as a leaf. A node blowup could create coefficient
one only from an edge $(1,1)$, which (C.CR.2)
excludes. Subsequent blowups preserve the valence
one of an existing coefficient-one vertex: at its
unique boundary node they merely subdivide that edge,
and at a smooth point its new exceptional coefficient
is zero.

Let $E_0$ be the connected exceptional fibre containing
$C$. In the boundary tree, the first nonexceptional
vertex on every path from $C$ meets $E_0$, so its image
passes through $p$. Its further image is therefore
a component of $\Gamma$ through $q$. By C.2.CR.2
there is only one such component, denoted $G$.

Suppose its generic transversal were a node.
Every strict conductor curve above $G$ then has
coefficient one. Thus each first nonexceptional
vertex just found is a leaf. Every strict conductor
curve is consequently such a leaf, attached at a
single point to $E_0$. No additional exceptional
tail can meet it, since that would raise its
boundary valence. The reduced conductor on $T$
is therefore a union of rational curves whose only
possible multibranch point is $p$; every component
passes through $p$, is smooth away from $p$, and
has only one branch at $p$.

All their images pass through $q$, where $\Gamma$
is smooth, so the conductor image has just the one
global component $G$. It also has no multibranch
point elsewhere. Indeed, at $x\ne q$ each point
of the normalization over $x$ receives just one
local branch of $G$. If $G$ had two distinct local
branches there, their receiving normalization
points would be distinct. The corresponding
irreducible hypersurface branches of $S$ must
intersect in a curve, because their distinct
local equations generate a height-two ideal in
the regular local ring of the threefold. That
intersection curve is a conductor-image branch
received by both normalization points, contradicting
their different singleton branch sets.

The normalization of $G$ is rational and all its
germs are unibranch. The normalization map is a
continuous bijection from a compact space to a
Hausdorff space and hence a homeomorphism.
Thus $\pi_1(G)=1$, contradicting (C.CR.1).
The generic cusp alternative is the only one left.
Its transverse normalization is unibranch with
conductor exponent two, so it contributes exactly
one strict coefficient-two curve, of degree one
over $G$.
\end{proof}


\begin{statement}{C.2.CR.4}{Proved}{A smooth critical curve with generic cusp is a constant cusp}
Let a reduced hypersurface germ $\{f=0\}$ in a smooth
complex threefold occur in a fibre whose nearby smooth
fibres are complex two-tori. Suppose that the full reduced
critical locus of $f$ near the point is a smooth curve
and that the generic transverse singularity is an ordinary
cusp. Then there are local holomorphic coordinates in which
\[
                            f=X^2+Y^3.             \tag{C.CR.15}
\]
An analytic unit and a common full-fibre multiplicity may
be retained in the original map: they do not change this
assertion about the reduced hypersurface.
\end{statement}
\begin{proof}
Choose coordinates $(x,y,z)$ with critical curve
$\{x=y=0\}$. The rank-one quadratic form along this curve
can be written
\[
 a(z)x^2.
\]
To obtain these coordinates without a meromorphic change,
divide the matrix of its coefficients by its common
vanishing order at $z=0$. A primitive symmetric rank-one
matrix has a unit entry, factors as a unit times the
square of a primitive linear form, and that linear form
extends to an invertible holomorphic change of the normal
coordinates. The intrinsic cubic coefficient on its
kernel is then $b(z)y^3$. The generic cusp hypothesis says
that both $a$ and $b$ are nonzero functions; shrink the
disc so that their only possible zero is at $z=0$.

If $a(0)\ne0$, the Morse lemma with parameters gives
\[
 f=X^2+y^3u(y,z),\qquad u(0,z)=b(z)
\]
up to multiplication of $b$ by a holomorphic unit.
If $u(0,0)=0$, the reduced support of $y^3u(y,z)$ is
singular: $u$ has a component through the origin
different from $y=0$, because $u(0,z)$ is not identically
zero. The arbitrary-suspension clause of A.5.17 then
puts a sphere of nonzero self-intersection into a
nearby actual torus fibre, which is impossible.
Thus $u$ is a unit, and replacing $y$ by a holomorphic
cube root of $u$ times $y$ gives (C.CR.15).

It remains to exclude $a(0)=0$. Consider the actual
holomorphic perturbation
\[
                  f_{\epsilon,\delta}
                      =f+\epsilon x^2+\delta y^3.
 \tag{C.CR.16}
\]
Choose small parameters so that $a(z)+\epsilon$ has a
simple zero $z_\epsilon$ and
$b(z_\epsilon)+\delta\ne0$. Such choices exist:
regular values of the nonconstant function $a$ are
dense near zero, and finitely many values of $\delta$
are then excluded.

We specify the germ produced at
$(0,0,z_\epsilon)$. Weierstrass preparation in $y$,
followed by a Tschirnhaus substitution, writes it as
a holomorphic unit times
\[
             Y^3+x^2\bigl(\beta(x,z)Y+\gamma(x,z)\bigr).
 \tag{C.CR.17}
\]
Here the absence of a linear transverse jet gives
the divisibility of the constant term by $x^2$,
and the rank-one quadratic form gives the divisibility
of the coefficient of $Y$ by $x^2$. More explicitly,
before the Tschirnhaus substitution the coefficient
of $y^2$ is divisible by $x$, while those of $y$
and of $1$ are divisible by $x^2$; these divisibilities
are preserved by that substitution.
The function $\gamma(0,z)$ is a unit multiple of
$a(z)+\epsilon$. Consequently
\[
             Z=\beta(x,z)Y+\gamma(x,z)
\]
is an invertible change of the curve coordinate at
the chosen simple zero. Equation (C.CR.17) is thus
a unit times $Y^3+x^2Z$. This remaining unit can
also be removed: this polynomial is homogeneous
of degree three, so scaling all three coordinates
by a holomorphic cube root of the unit gives the
exact equation
\[
                            Y^3+x^2Z.
 \tag{C.CR.18}
\]
A.5.15 supplies a compact spherical cycle of square $-4$
in an arbitrarily small nonzero level of this germ.
Indeed its Hirzebruch formula is
$-2(b/\gcd(b,c))(c/\gcd(b,c))$, here with
$(b,c)=(2,1)$. The assertion is about the homology
class of the explicit continuous map from $S^2$;
no embedded representative is needed.

For completeness we prove the boundary comparison
needed to put this spherical cycle in the original fibre.
Choose a sufficiently small closed ball whose
boundary meets the smooth critical curve transversely
and is transverse to the smooth stratum of $f=0$.
On the compact link of the critical curve, both
$a$ and $b$ are units. For all sufficiently small
$(\epsilon,\delta)$ the parametrized Morse and
cusp-coordinate changes therefore give holomorphic
product-cusp charts there, retaining $z$ as curve
coordinate. In each such chart the vector
$\partial/\partial z$ is tangent to every level
of $f_{\epsilon,\delta}$. Along the critical link,
an appropriate real radial vector in this direction
is transverse to the fixed sphere. This remains
true on a neighbourhood of that link, uniformly
for the small perturbations. Finitely many charts
suffice. On the rest of the boundary ordinary
compactness and transversality at parameter zero
give the same conclusion for small target values.
Thus all sufficiently small nonzero levels have
uniformly transverse boundary. The full-critical-
locus hypothesis ensures that no other boundary
singular curve has been omitted from this argument.

For fixed small parameters, the nonzero critical
values inside the ball are finite. The local
spherical cycle from (C.CR.18) can therefore be transported
through regular target values to a fixed small
value $t_1$, avoiding those critical values.
Choose the parameters small enough that every
critical value throughout the segment joining
them to zero has modulus less than $|t_1|/2$.
This is possible by compactness: every limit of
critical points is a critical point of $f$, and
its value is zero; the just-proved product charts
control the boundary as well. The family of
$t_1$-levels along this parameter segment is a
proper submersion with transverse boundary.
Ehresmann transport identifies it with the
original $f=t_1$ level and preserves the square
of the transported spherical class.

The original level lies in a smooth nearby torus
fibre. A two-torus is aspherical, so every sphere
in it is null-homotopic and has zero homology
class and self-intersection. This contradicts
the nonzero square obtained above. Hence $a(0)\ne0$,
and the first part of the proof gives (C.CR.15).
If the original map is a unit times $f^m$, absorb
the unit by its local $m$th root before this
comparison; the resulting sphere lies in the
actual level with parameter $t_1^m$.
\end{proof}

\begin{statement}{C.2.CR.5}{Proved}{A contracted horizontal section is elliptic}
In the situation of C.2.CR.1 one has
\[
                              g(B)=1.              \tag{C.CR.19}
\]
In particular the remaining exact Euler expression is
\[
                            e(S)=n-v+N_0.          \tag{C.CR.20}
\]
\end{statement}
\begin{proof}
By C.2.CR.2 the only other possibility is $g=2$.
The conductor image is then smooth and reduced at
$q=h(C)$, and its generic transversal is a node or
a cusp. C.2.CR.3 excludes the node. For the cusp,
the reduced critical locus near $q$ is precisely
that smooth conductor curve: every positive-dimensional
singular locus of a reduced hypersurface surface
belongs to its nonnormal locus, and there are no
other isolated points after the neighbourhood is
shrunk. C.2.CR.4 makes the germ a product cusp.
Its normalization is smooth, contrary to the
exceptional genus-two curve contracted to the
normalization point $p$. Thus $g=2$ is excluded.
Substitution in (C.CR.3) gives (C.CR.20).
\end{proof}


\begin{statement}{C.2.CR.6}{Proved}{The elliptic contraction has a planar conductor and only double transverse branches at its image}
In the remaining case $g=1$ of C.2.CR.1, let
$p=\pi(C)$ and $q=\eta(p)$. The conductor-image
scheme is a local complete intersection at $q$.
More precisely, ambient coordinates can be chosen
so that its ideal is
\[
                         (x,g(y,z)).               \tag{C.CR.21}
\]
Every global conductor component passing through
$q$ has generic transverse type $A_k$ for some
$k\ge1$. For type $A_{2s-1}$, its two transverse
normalization branches have conductor coefficient
$s$ each; their global covering degrees sum to two.
For type $A_{2s}$ there is one coefficient-$2s$
normalization branch, of covering degree one.
\end{statement}
\begin{proof}
The residue comparison (C.CR.12) is valid also
when $g=1$. By (C.CR.13) its nonzero one-dimensional
quotient requires an element of the conductor
summand $I_p$ with nonzero ambient linear part.
Call this element $x$ and extend it to coordinates
$(x,y,z)$ on the smooth threefold.

The full conductor ideal contains $x$ and has a
pure one-dimensional Cohen--Macaulay quotient.
After quotienting the ambient local ring by $x$,
its remaining ideal is an unmixed height-one
ideal in the regular two-dimensional ring
$\mathbb C\{y,z\}$. Such an ideal is principal:
it is the intersection of its height-one primary
localizations, and each is a power of a principal
prime in this factorial ring. Equivalently its
quotient has depth one, so the ideal is a rank-one
reflexive module on this smooth surface and
is invertible. This proves (C.CR.21), retaining
all multiplicities in $g$.

We give the precise adjunction comparison behind
the transverse multiplicity assertion. For a
reduced plane curve germ $D\subset W$, let
$\nu:\widetilde D\to D$ be its normalization,
and let $\rho:\widetilde W\to W$ be an embedded
resolution with smooth disjoint strict transform
$\widehat D$. Finite curve duality identifies
\[
 \mathfrak c_D\,\omega_D=\nu_*\omega_{\widetilde D}.
\]
The residue sequence on $\widetilde W$ and
Grauert--Riemenschneider vanishing identify the
right side with the residue of
$\rho_*\omega_{\widetilde W}(\widehat D)$.
In a local generator $du\wedge dv/f_D$ of the
ambient meromorphic residue description, a
conductor element $a$ therefore lifts to a form
\[
                     a\,du\wedge dv/f_D
\]
whose pullback has no exceptional pole. On the
first plane blowup, with exceptional curve $E$,
the canonical divisor has coefficient one and
the total transform of $D$ has coefficient
$m=\operatorname{mult}_0 D$. The coefficient
of this pulled-back form along $E$ is
$\operatorname{ord}_0(a)+1-m$. Hence
\[
 \mathfrak c_D\subset\mathfrak m_D^{\,m-1}
 \quad\hbox{in the sense that every ambient
 representative has order at least }m-1.
 \tag{C.CR.22}
\]
Adding a multiple of $f_D$ does not change this
order bound, since $f_D$ has order $m$.

At a general smooth point of any conductor
component through $q$, the ambient coordinate
$x$ remains in its conductor ideal and has
a nonzero transverse linear part: its zero
set is a smooth ambient plane containing the
conductor curve. Applying (C.CR.22) to a
transverse plane slice gives $m\le2$.
The slice is singular, so $m=2$.
Weierstrass preparation, completion of the
square, and absorption of a unit then put
a reduced multiplicity-two plane curve in
the form $u^2-v^{k+1}=0$.

For $k=2s-1$, the two branches are
$u=\pm v^s$; their intersection multiplicity
is $s$, and the conductor on each normalized
branch is $(v^s)$. For $k=2s$, the normalization
is $v=t^2$, $u=t^{2s+1}$, whose semigroup
has conductor $2s$, giving $(t^{2s})$.
These calculations specify both the coefficients
and the numbers of transverse branches; finite
normalization over the global image component
expresses those numbers as the asserted sums
of covering degrees.
\end{proof}




